% ============================================================
% M3 S3 练习件·W3 臂H 样板件 —— 衔接节（2.8 课前衔接·升格全制式·练习位独立）＝练习件线 S3 线末位片
% 生成：M3 成卷轮 S3 W3 臂H 量产代理｜2026-09-14｜编译 cwd＝本目录（中文路径禁 \input 绝对路径）
%   xelatex main-true.tex ×2 → 含详解印本；xelatex main-false.tex ×2 → 纯题版（单源双本）
% 输入链（全只读）：题面库 衔接节.md（题面逐字装配）＋ 衔接节-答案侧.md（值逐字回嵌）
%   ＋ 定稿/批E-衔接节.md（详解回嵌源，§七 补产全文区）；toolchain＝qp-m3.sty（钉后版
%   md5 7c3930362be8a0a2bdf21bbf8ac16573，本地挂载副本，破损即停）。
% 母版体例：照抄《课时01-坐标法/练习件母版体例说明.md》骨架与门谱；件标题位照 \xjkeshi
%   口径（sty 要件⑤，「衔接」前缀恒带，导学件衔接节同款）；知识点标照导学件知识导学
%   分册（一判别式／二韦达／三方程组解法）＋探究点域（四＝位置关系：判定/弦长/切线）。
% 键账：件内 ansblock 键集 ≡ 题面库片练键集 ≡ manifest 练键序 练1~16（16 键，零缺漏零多余）；
%   拓区隔离：2章-拓/测/滚 键不入本件（S4 拓展册收）；导学键 G1~G5 与探究键 探1~探6
%   属导学件域，亦不入（练习位独立）。
% 卷式例外案B：练习件不属卷式——答案块物理落点恒为题后 ansblock，禁卷末附卷制；
%   全线括线模（导言 \ansblockgrayfalse 一次置定，件内不切换；灰底盒不可跨栏断，练习件禁用）。
% 题序＝manifest 键序 练1~16（零重排）；印面题号 1~16 连号（\tihao 号＝\ansitem 号同键）；
%   三档切位 1~7／8~14／15~16（照库序切位，非难度重排；难度读数 简8＋中6＋难2 照批E §四）。
% 槽配实录：单选2（练14/15）＋多选0（manifest 期望值在案）＋填空1（练2）＋解答13（其余）；
%   衔接节练习位为存量13保真＋补3 的预备课位槽配，非 10选4填2解 基线（如实登记）。
% 选项槽位：默认四连排（0.25\linewidth−0.25\lxhang），超宽降档梯 四连排→两连排
%   （0.5\linewidth−0.5\lxhang−0.25em）→单列 \lxopt；禁缩字号/负 kern/删标点凑宽。
%   本片实测降档：题14/题15 四连排严格档全旗（C/D 槽 60.4pt、四槽 64.1~64.6pt＞预警线
%   46.2pt）→ 双双落两连排（52%~61% 全过）；读数见值台账「降档登记」。
% 解答书写区：13 题按小问数定高——练1（4问）/练8（3问）32mm；练4/5/6（2问）24mm；
%   无小问 9 题 16mm 默认档（cap 40mm 内；纯题档吞 ansblock 不吞 \liubai）。
% 填空印答：练2 空位 \kongda{21}（详解版印答、纯题版退定宽空线，两档等形）。
% ============================================================
\documentclass[fontset=none]{ctexart}
\usepackage{qp-m3}
% —— 印面逐字符号路由补钉（承导学件母版；− 走 TNR、▱ NSC 子块；禁改 sty 保 md5） ——
\xeCJKDeclareCharClass{Default}{"2212}%
\setCJKfamilyfont{nscblk}[Path=C:/提示词/工作区/字替对照-0909/variantF/fonts/,BoldFont=NSC-w600.ttf]{NSC-w500.ttf}%
\xeCJKDeclareSubCJKBlock{nscblk}{"25B1}%
\newcommand{\pxparallelogram}{{\CJKfamily{nscblk}▱}}%
% —— 断行微调（承导学件母版实测档，三零门承重；\emergencystretch 不另设） ——
\xeCJKsetup{CJKglue={\hskip 0.10em plus 0.08em minus 0.05em}}%
% —— 练习件全线括线模（S3 波次方案 §四.3：导言一次置定、件内不切换） ——
\ansblockgrayfalse
% —— 页脚件名（qp-headfoot 复用入口） ——
\renewcommand{\qpzhangming}{第二章\quad 平面解析几何}%
\renewcommand{\qpjianming}{练习件}%

\begin{document}
%装配轮S6三本跨本连续页码（G7方案B）setcounter 注入位
\setcounter{page}{76}

\xjkeshi{2.8 直线与圆锥曲线的位置关系}
\begin{multicols}{2}
\emergencystretch=1em

% ============================================================
% 基础巩固（题1~7＝练1~7：解答6＋填空1；题后紧跟答案制，全线括线 ansblock）
% ============================================================
\dangbadge{基}{础}{巩}{固}

\tihao{1}\zu{韦达对称式四小问×解答} \tieside{简单(知识点二)}若\(x_1\)、\(x_2\)是方程\(2x^{2}-4x+1=0\)的两根，不解方程，求下列各式的值：
\bindp (1) \(x_1^{2}+x_2^{2}\)；(2) \(|x_1-x_2|\)；(3) \(x_1/x_2\)（\(x_1>x_2\)）；(4) \(1/x_1-x_2^{2}\)（\(x_1>x_2\)）．
\liubai[32mm]{解答书写区}

\begin{ansblock}[2章-练-衔接-1]
% ans:2章-练-衔接-1
\ansitem{1}{(1) 3；(2) √2；(3) 3+2√2；(4) 1/2}
\ansnote{详解}{韦达定理：\(x_1+x_2=2\)、\(x_1x_2=1/2\)．(1) \(x_1^{2}+x_2^{2}=(x_1+x_2)^{2}-2x_1x_2=4-1=3\)．\par\noindent (2) \(|x_1-x_2|=\sqrt{(x_1+x_2)^{2}-4x_1x_2}=\sqrt{2}\)．\par\noindent (3) 由求根公式定序：\(x_1=(2+\sqrt{2})/2\)、\(x_2=(2-\sqrt{2})/2\)，故\(x_1/x_2=(2+\sqrt{2})/(2-\sqrt{2})=(2+\sqrt{2})^{2}/2=3+2\sqrt{2}\)．\par\noindent (4) 由\(x_1x_2=1/2\)得\(1/x_1=2x_2\)；又\(x_2\)是原方程的根，\(x_2^{2}=2x_2-1/2\)，故\(1/x_1-x_2^{2}=2x_2-(2x_2-1/2)=1/2\)．}
\end{ansblock}

\tihao{2}\zu{高次对称式×填空} \tieside{简单(知识点二)}已知\(x_1\)、\(x_2\)是方程\(x^{2}-x-1=0\)的两实根，则\(2x_1^{5}+5x_2^{3}=\)\kongda{21}．

\begin{ansblock}[2章-练-衔接-2]
% ans:2章-练-衔接-2
\ansitem{2}{21}
\ansnote{详解}{由方程得\(x^{2}=x+1\)，即\(x_1^{2}=x_1+1\)、\(x_2^{2}=x_2+1\)，且\(x_1+x_2=1\)．\(2x_1^{5}=2x_1(x_1^{2})^{2}=2x_1(x_1+1)^{2}=2x_1(3x_1+2)=6x_1^{2}+4x_1=10x_1+6\)；\(5x_2^{3}=5x_2(x_2+1)=5x_2^{2}+5x_2=10x_2+5\)．合计\(10(x_1+x_2)+11=10+11=21\)．}
\end{ansblock}

\tihao{3}\zu{根与系数互求×解答} \tieside{简单(知识点二)}已知\(2+\sqrt{3}\)是关于\(x\)的方程\(x^{2}-4x+c=0\)的一个根，求方程的另一根及\(c\)的值．
\liubai[16mm]{解答书写区}

\begin{ansblock}[2章-练-衔接-3]
% ans:2章-练-衔接-3
\ansitem{3}{另一根 2−√3；c=1}
\ansnote{详解}{两根之和为4，故另一根\(x_2=4-(2+\sqrt{3})=2-\sqrt{3}\)；两根之积为\(c\)，故\(c=(2+\sqrt{3})(2-\sqrt{3})=4-3=1\)．双路验算：把\(2+\sqrt{3}\)代入方程得\(7+4\sqrt{3}-8-4\sqrt{3}+c=0\)，亦得\(c=1\)．}
\end{ansblock}

\tihao{4}\zu{含参两问×解答} \tieside{中档(知识点二)}关于\(x\)的二次方程\(kx^{2}+2x-ck=0\)（\(c>0\)，\(c\neq 1\)）有两个不等的实数根，其中一个是\(x=-c\)．
\bindp (1) 求另一个根；(2) 写出\(k\)关于\(c\)的表达式，并求\(k\)的取值范围．
\liubai[24mm]{解答书写区}

\begin{ansblock}[2章-练-衔接-4]
% ans:2章-练-衔接-4
\ansitem{4}{(1) 1；(2) k=2/(c−1)，k>0 或 k<−2}
\ansnote{详解}{(1) 两根之积\(=-ck/k=-c\)，故另一根\((-c)/(-c)=1\)．(2) 两根之和\(1-c=-2/k\)，故\(k=2/(c-1)\)．由\(c>0\)且\(c\neq 1\)：\(c>1\)时\(c-1>0\)，\(k>0\)；\(0<c<1\)时\(-1<c-1<0\)，\(k=2/(c-1)<2/(-1)=-2\)，即\(k<-2\)．故\(k>0\)或\(k<-2\)．（另“两不等实根”前提恒相容：\(\Delta=4+4ck^{2}>0\)恒成立．）}
\end{ansblock}

\tihao{5}\zu{恒有根证明＋求参×解答} \tieside{中档(知识点一)}已知关于\(x\)的一元二次方程\(x^{2}+(2m+1)x+m-2=0\)．
\bindp (1) 求证：无论\(m\)取何值，此方程总有两个不相等的实数根；
\bindp (2) 若方程有两个实数根\(x_1\)、\(x_2\)，且\(x_1+x_2+3x_1x_2=1\)，求\(m\)的值．
\liubai[24mm]{解答书写区}

\begin{ansblock}[2章-练-衔接-5]
% ans:2章-练-衔接-5
\ansitem{5}{(1) 证毕（Δ=4m²+9>0）；(2) m=8}
\ansnote{详解}{(1) \(\Delta=(2m+1)^{2}-4(m-2)=4m^{2}+4m+1-4m+8=4m^{2}+9\geqslant 9>0\)，故无论\(m\)取何值，方程总有两个不相等的实数根．\par\noindent (2) 韦达定理：\(x_1+x_2=-(2m+1)\)、\(x_1x_2=m-2\)，代入\(x_1+x_2+3x_1x_2=1\)得\(-(2m+1)+3(m-2)=1\)，即\(m-7=1\)，解得\(m=8\)．}
\end{ansblock}

\tihao{6}\zu{存在性＋整除枚举×解答} \tieside{难(知识点一)}已知\(x_1\)、\(x_2\)是一元二次方程\(4kx^{2}-4kx+k+1=0\)的两个实数根．
\bindp (1) 是否存在实数\(k\)，使得\((2x_1-x_2)(x_1-2x_2)=-3/2\)成立？若存在，求出\(k\)的值；若不存在，请说明理由；
\bindp (2) 求使\(x_1/x_2+x_2/x_1-2\)的值为整数的实数\(k\)的整数值．
\liubai[24mm]{解答书写区}

\begin{ansblock}[2章-练-衔接-6]
% ans:2章-练-衔接-6
\ansitem{6}{(1) 不存在；(2) k=−2，−3，−5}
\ansnote{详解}{两实根须\(4k\neq 0\)且\(\Delta=(-4k)^{2}-4\cdot 4k\cdot(k+1)=-16k\geqslant 0\)，故\(k<0\)；韦达定理：\(x_1+x_2=1\)、\(x_1x_2=(k+1)/(4k)\)．\par\noindent (1) \((2x_1-x_2)(x_1-2x_2)=2x_1^{2}+2x_2^{2}-5x_1x_2=2[(x_1+x_2)^{2}-2x_1x_2]-5x_1x_2=2-9x_1x_2=2-9(k+1)/(4k)\)．令其等于\(-3/2\)：\(18(k+1)=28k\)，得\(10k=18\)，\(k=9/5>0\)，与实根前提\(k<0\)矛盾，故不存在．\par\noindent (2) \(x_1/x_2+x_2/x_1-2=(x_1+x_2)^{2}/(x_1x_2)-4=4k/(k+1)-4=-4/(k+1)\)．\(k\)为负整数且\(k\neq -1\)，需\(k+1\)整除4，\(k+1\in\{-1,-2,-4\}\)，即\(k\in\{-2,-3,-5\}\)．}
\end{ansblock}

\tihao{7}\zu{勾股＋韦达×解答} \tieside{难(知识点二)}已知Rt\(\triangle ABC\)中，两直角边长为方程\(x^{2}-(2m+7)x+4m(m-2)=0\)的两根，且斜边长为13，求\(S_{\triangle ABC}\)的值．
\liubai[16mm]{解答书写区}

\begin{ansblock}[2章-练-衔接-7]
% ans:2章-练-衔接-7
\ansitem{7}{30}
\ansnote{详解}{设两直角边为\(a\)、\(b\)：\(a+b=2m+7\)、\(ab=4m(m-2)\)．勾股定理：\(169=a^{2}+b^{2}=(a+b)^{2}-2ab=(2m+7)^{2}-8m(m-2)\)，展开得\(4m^{2}+28m+49-8m^{2}+16m=169\)，整理得\(m^{2}-11m+30=0\)，解得\(m=5\)或\(m=6\)．\par\noindent 实根检验：\(\Delta=(2m+7)^{2}-16m(m-2)\)，\(m=6\)时\(\Delta=361-384=-23<0\)，舍；\(m=5\)时\(\Delta=289-240=49>0\)，且\(a+b=17>0\)、\(ab=60>0\)均正，符合．故\(S_{\triangle ABC}=ab/2=30\)．}
\end{ansblock}

% ============================================================
% 综合提升（题8~14＝练8~14：解答6＋单选1，方程组域收束＋位置关系首题）
% ============================================================
\dangbadge{综}{合}{提}{升}

\tihao{8}\zu{解三方程组×解答} \tieside{简单(知识点三)}解下列方程组：
\bindp (1) \(\{y^{2}=3x+1,\ 3x-y=1\}\)；(2) \(\{(x+y)^{2}-2(x+y)-3=0,\ x-y=-1\}\)；(3) \(\{9x^{2}-4y^{2}-32=0,\ 3x-2y-4=0\}\)．
\liubai[32mm]{解答书写区}

\begin{ansblock}[2章-练-衔接-8]
% ans:2章-练-衔接-8
\ansitem{8}{(1) x=0,y=−1 或 x=1,y=2；(2) x=1,y=2 或 x=−1,y=0；(3) x=2,y=1}
\ansnote{详解}{(1) 由\(3x-y=1\)得\(y=3x-1\)，代入得\(9x^{2}-6x+1=3x+1\)，即\(9x^{2}-9x=0\)，解得\(x=0\)或\(x=1\)，回代得\(y=-1\)或\(y=2\)．\par\noindent (2) 整体设元\(t=x+y\)：\(t^{2}-2t-3=0\)，得\(t=3\)或\(t=-1\)；配\(x-y=-1\)：\(t=3\)时\(x=1\)、\(y=2\)；\(t=-1\)时\(x=-1\)、\(y=0\)．\par\noindent (3) 平方差：\(9x^{2}-4y^{2}=(3x-2y)(3x+2y)=32\)，而\(3x-2y=4\)，故\(3x+2y=8\)；与\(3x-2y=4\)联立得\(x=2\)、\(y=1\)．（注意：消元所得根必须回代求另一未知数，防丢解．）}
\end{ansblock}

\tihao{9}\zu{分式方程组×解答} \tieside{简单(知识点三)}解方程组\(\{1/x+1/y=0,\ xy=-4\}\)．
\liubai[16mm]{解答书写区}

\begin{ansblock}[2章-练-衔接-9]
% ans:2章-练-衔接-9
\ansitem{9}{x=−2,y=2 或 x=2,y=−2}
\ansnote{详解}{\(xy\neq 0\)，第一式同乘\(xy\)：\(x+y=0\)，即\(x=-y\)．代入\(xy=-4\)：\(-y^{2}=-4\)，\(y=\pm 2\)，故\((x,y)=(-2,2)\)或\((2,-2)\)．}
\end{ansblock}

\tihao{10}\zu{换元×解答} \tieside{中档(知识点三)}解方程组\(\{x-y-3\sqrt{x-y}=10,\ xy=-6\}\)．
\liubai[16mm]{解答书写区}

\begin{ansblock}[2章-练-衔接-10]
% ans:2章-练-衔接-10
\begingroup\rightskip=0pt plus 1fil\relax
\ansitem{10}{x=(25+√601)/2,y=(−25+√601)/2 或 x=(25−√601)/2,y=(−25−√601)/2}
\endgroup
\ansnote{详解}{设\(A=\sqrt{x-y}\geqslant 0\)：\(A^{2}-3A-10=0\)，\((A-5)(A+2)=0\)，得\(A=5\)（负根舍），故\(x-y=25\)．配\(xy=-6\)：\(x(x-25)=-6\)，即\(x^{2}-25x+6=0\)，解得\(x=(25\pm\sqrt{625-24})/2=(25\pm\sqrt{601})/2\)；\(y=x-25=(-25\pm\sqrt{601})/2\)（\(\pm\)号对应取同号）．}
\end{ansblock}

\tihao{11}\zu{对称构造×解答} \tieside{简单(知识点三)}解方程组\(\{x+y=5,\ xy=6\}\)．
\liubai[16mm]{解答书写区}

\begin{ansblock}[2章-练-衔接-11]
% ans:2章-练-衔接-11
\ansitem{11}{x=2,y=3 或 x=3,y=2}
\ansnote{详解}{法一（代入）：\(y=5-x\)代入\(xy=6\)得\(x^{2}-5x+6=0\)，\(x=2\)或\(x=3\)，对应\(y=3\)或\(y=2\)．法二（对称构造）：\(x\)、\(y\)是\(t^{2}-5t+6=0\)的两根，同解．}
\end{ansblock}

\tihao{12}\zu{方程组两解定参×解答} \tieside{中档(知识点三)}当\(m\)为何值时，关于\(x\)、\(y\)的方程组\(\{x^{2}/4+y^{2}=1,\ y=x+m\}\)有两个解．
\liubai[16mm]{解答书写区}

\begin{ansblock}[2章-练-衔接-12]
% ans:2章-练-衔接-12
\ansitem{12}{−√5<m<√5}
\ansnote{详解}{把\(y=x+m\)代入\(x^{2}/4+y^{2}=1\)：\(x^{2}/4+(x+m)^{2}=1\)，整理得\(5x^{2}+8mx+4m^{2}-4=0\)．方程组有两解\(\iff\)该二次方程有两不等实根\(\iff\Delta=64m^{2}-20(4m^{2}-4)=80-16m^{2}>0\)，即\(m^{2}<5\)，故\(-\sqrt{5}<m<\sqrt{5}\)．}
\end{ansblock}

\tihao{13}\zu{最值×解答} \tieside{中档(知识点三)}实数\(x\)、\(y\)满足\(x^{2}/4+y^{2}=1\)，求\(y-x\)的最大值．
\liubai[16mm]{解答书写区}

\begin{ansblock}[2章-练-衔接-13]
% ans:2章-练-衔接-13
\ansitem{13}{√5（此时 x=−4√5/5、y=√5/5）}
\ansnote{详解}{设\(m=y-x\)，即\(y=x+m\)，代入\(x^{2}/4+y^{2}=1\)：\(5x^{2}+8mx+4(m^{2}-1)=0\)．实数解存在\(\iff\Delta=64m^{2}-80(m^{2}-1)\geqslant 0\)，即\(m^{2}\leqslant 5\)，故\(m\)的最大值为\(\sqrt{5}\)．取等时关于\(x\)的方程为\(5x^{2}+8\sqrt{5}x+16=0\)，解得\(x=-8\sqrt{5}/10=-4\sqrt{5}/5\)，\(y=x+\sqrt{5}=\sqrt{5}/5\)，点在椭圆上，符合．}
\end{ansblock}

\tihao{14}\zu{等轴双曲线弦长逆求×单选} \tieside{中档(知识点四)}已知等轴双曲线的中心在原点，焦点在\(x\)轴上，与直线\(y=(1/2)x\)交于\(A\)、\(B\)两点，若\(|AB|=2\sqrt{15}\)，则该双曲线的方程为\nobreak(\kongwei)
\bindopt {\fontsize{10.5pt}{\qpTJgLeadLiti}\selectfont \makebox[\dimexpr0.5\linewidth-0.5\lxhang-0.25em\relax][l]{A．\(x^{2}-y^{2}=6\)}\makebox[\dimexpr0.5\linewidth-0.5\lxhang-0.25em\relax][l]{B．\(x^{2}-y^{2}=9\)}\par}
\bindopt {\fontsize{10.5pt}{\qpTJgLeadLiti}\selectfont \makebox[\dimexpr0.5\linewidth-0.5\lxhang-0.25em\relax][l]{C．\(x^{2}-y^{2}=16\)}\makebox[\dimexpr0.5\linewidth-0.5\lxhang-0.25em\relax][l]{D．\(x^{2}-y^{2}=25\)}\par}

\begin{ansblock}[2章-练-衔接-14]
% ans:2章-练-衔接-14
\ansitem{14}{B（x²−y²=9）}
\ansnote{详解}{设双曲线方程为\(x^{2}-y^{2}=\lambda\)（\(\lambda>0\)，焦点在\(x\)轴的等轴双曲线）．与\(y=x/2\)联立：\((3/4)x^{2}=\lambda\)，\(x=\pm 2\sqrt{\lambda/3}\)，\(|\Delta x|=4\sqrt{\lambda/3}\)；\(|AB|=\sqrt{1+k^{2}}\,|\Delta x|=(\sqrt{5}/2)\cdot 4\sqrt{\lambda/3}=2\sqrt{5\lambda/3}=2\sqrt{15}\)，得\(5\lambda/3=15\)，\(\lambda=9\)，方程为\(x^{2}-y^{2}=9\)，选B．}
\end{ansblock}

% ============================================================
% 思维探索（题15~16＝练15~16：单选1＋解答1，0913补题入槽，居末档照库序）
% ============================================================
\dangbadge{思}{维}{探}{索}

\tihao{15}\zu{平行切线×单选} \tieside{简单(知识点四)}与直线\(2x-y+4=0\)平行的抛物线\(y=x^{2}\)的切线方程为\nobreak(\kongwei)
\bindopt {\fontsize{10.5pt}{\qpTJgLeadLiti}\selectfont \makebox[\dimexpr0.5\linewidth-0.5\lxhang-0.25em\relax][l]{A．\(2x-y+3=0\)}\makebox[\dimexpr0.5\linewidth-0.5\lxhang-0.25em\relax][l]{B．\(2x-y-3=0\)}\par}
\bindopt {\fontsize{10.5pt}{\qpTJgLeadLiti}\selectfont \makebox[\dimexpr0.5\linewidth-0.5\lxhang-0.25em\relax][l]{C．\(2x-y+1=0\)}\makebox[\dimexpr0.5\linewidth-0.5\lxhang-0.25em\relax][l]{D．\(2x-y-1=0\)}\par}

\begin{ansblock}[2章-练-衔接-15]
% ans:2章-练-衔接-15
\ansitem{15}{D（2x−y−1=0）}
\ansnote{详解}{设切线为\(2x-y+m=0\)，即\(y=2x+m\)，代入\(y=x^{2}\)：\(x^{2}-2x-m=0\)．相切\(\iff\Delta=4+4m=0\)，得\(m=-1\)，切线为\(2x-y-1=0\)（切点\((1,1)\)），选D．}
\end{ansblock}

\tihao{16}\zu{开放条件判型×解答} \tieside{简单(知识点四)}若关于\(x\)、\(y\)的方程\(x^{2}/(3-t)+y^{2}/(t-1)=1\)表示的是曲线\(C\)，给出下列三个条件：①曲线\(C\)是椭圆，②长轴在\(y\)轴上，③长轴在\(x\)轴上．请选择其中2个条件与已知组成命题，并求出\(t\)的取值范围．
\liubai[16mm]{解答书写区}

\begin{ansblock}[2章-练-衔接-16]
% ans:2章-练-衔接-16
\ansitem{16}{选①②：2<t<3；选①③：1<t<2}
\ansnote{详解}{条件①表示椭圆\(\iff\)两分母同正且不等：\(3-t>0\)、\(t-1>0\)、\(3-t\neq t-1\)，即\(1<t<3\)且\(t\neq 2\)．选①②（长轴在\(y\)轴）：\(t-1>3-t>0\)，解得\(2<t<3\)；选①③（长轴在\(x\)轴）：\(3-t>t-1>0\)，解得\(1<t<2\)．}
\end{ansblock}

% ---- 尾块（\tailfill 置 \end{multicols} 前；无参＝内容尾随框，每件恰 1 处） ----
\tailfill

\end{multicols}
\end{document}
